\documentclass[../main.tex]{subfiles}

\begin{document}

\chapter{Infrastruttura di simulazione per SAND}
  \label{ch:infrastruttura-di-simulazione-per-sand}
  \begin{figure}[ht]
    \centering
    \includegraphics[width=.8\textwidth]{sim-infrastructure/sw-chain}
    \caption{
      Schema della catena software di simulazione e ricostruzione di SAND
    }
    \label{fig:sw-chain}
  \end{figure}
  L'infrastruttura software di simulazione per SAND, schematizzata in
  figura~\ref{fig:sw-chain} è composta da quattro elementi principali:
  \begin{itemize}
    \item le interazioni dei neutrini con la materia sono generate da GENIE;
    \item le particelle scaturite dalle interazioni dei neutrini sono
    propagate da EDep-sim (un wrapper attorno a GEANT4);
    \item la digitalizzazione e la ricostruzione del segnale sono
    svolte da SAND-Reco;
    \item i risultati della ricostruzione di SAND, aggregati a quelli degli
    altri componenti del ND, sono memorizzati su dei file in
    formato \textit{Common Analysis File} (CAF) dal ND\_CAFMaker
    (non rappresentato in figura~\ref{fig:sw-chain}).
  \end{itemize}
  In questo capitolo saranno presentati questi quattro strumenti, con un
  particolare focus su sandreco.

  \subfile{genie.tex}

  \subfile{edepsim.tex}

  \subfile{sandreco/sandreco.tex}

  \subfile{caf.tex}

\end{document}